\section{Methods}

\subsection{Frozen EEG features}
We used \texttt{brain-bzh/reve-base} with the encoder frozen and in evaluation
mode. HBN and MIPDB recordings were mapped to the ordered EGI HydroCel E1--E128
layout without spatial interpolation. Signals were resampled to 200 Hz,
band-pass filtered at 0.5--99.5 Hz, standardized and clipped at 15 standard
units, and divided into non-overlapping two-second windows. At most 120 seconds
were retained per participant. For MIPDB, standardization was fitted separately
to each selected recording before windowing, without age labels. Windows did
not cross condition boundaries. They came from the eyes-open and eyes-closed
segments of block01, selected in acquisition order. This is recording-level
normalization, not a global scaler fitted to the HBN training set.

Features from the last two encoder layers were cached once for the primary
comparison. The layer-wise follow-up also extracted the two preceding layers.
Each head predicted age for individual windows; the arithmetic mean of those
predictions gave one age estimate per participant.

\subsection{Prediction heads}
The baseline averages the final-layer tokens and applies an affine regressor.
The penultimate-layer probe applies the same readout to layer $-2$, so this
comparison changes depth without increasing the number of parameters.
The rich-statistics head adds a linear correction based on four summaries of
each token feature: standard deviation, range, mean absolute deviation, and
mean absolute value. These summaries allow a nonlinear dependence on the
original tokens, although the regression on the summaries is linear.

The multi-query head uses two trainable attention queries, initialized to
opposite signed basis vectors. Each query produces a weighted mean, standard
deviation, mean absolute deviation, and mean absolute value. Both queries also
receive the same unweighted token range. The correction uses the resulting ten
feature vectors. Both richer heads multiply their correction by 0.5 and
initialize its weights to zero. With matched seeds, they therefore start with
the same predictions as the baseline.

\begin{table}[tbp]
  \centering
  \small
  \caption{Primary heads. Feature dimension is 512. Parameter counts include
  all trainable head parameters; the encoder remains frozen.}
  \label{tab:heads}
  \begin{tabular}{p{0.26\linewidth}cp{0.40\linewidth}r}
  \toprule
  Head & Layer & Trainable components & Parameters \\
  \midrule
  Mean-pooled linear & $-1$ & Affine regression & \BaselineParameters{} \\
  Penultimate-layer linear & $-2$ & Affine regression & \LayerParameters{} \\
  Rich-statistics residual & $-1$ & Affine baseline and correction on four summaries & \RichParameters{} \\
  Multi-query rich-statistics & $-1$ & Affine baseline, two queries, and correction on ten summaries & \QueryParameters{} \\
  \bottomrule
  \end{tabular}
\end{table}

\subsection{Training and checkpoint selection}
For each head we trained seeds 33--42, shuffling windows with the corresponding
seed. Training used AdamW \citep{loshchilov2019adamw}, learning rate $10^{-4}$,
weight decay 0.05, batch size 64, a cosine OneCycleLR schedule
\citep{smith2018superconvergence}, gradient clipping at norm 1, and mean squared
error. The budget was 40 epochs with patience 7. We selected the earliest
checkpoint with the best participant-level HBN validation Pearson correlation.
For the primary comparison, neither external predictions nor age labels
influenced training or selection.

\subsection{Metrics}
For each seed, we computed Pearson correlation between predicted and observed
ages across external participants. The primary effect is the candidate's
correlation minus the matched baseline's correlation, averaged over the ten
seeds. We also report MAE and RMSE in years, $R^2$, and calibration. Calibration
fits observed age as intercept plus slope times predicted age; ideal values
are zero and one, respectively. These additional metrics describe prediction
quality but do not replace the primary Pearson endpoint.

\subsection{Uncertainty and decision rule}
We used 10,000 paired bootstrap iterations to estimate a percentile 95\%
confidence interval. Each iteration sampled ten seeds with replacement and
sampled external participants with replacement. The same participant sample
was used for every sampled seed and for both heads. We recomputed their
correlations and averaged the paired differences over sampled seeds. This
joint seed-and-participant bootstrap preserves the shared evaluation cohort;
participants are not independently resampled within each seed.

A separate one-sided sign-flip test enumerated all $2^{10}=1{,}024$ sign
assignments to the observed seed differences. Exactness requires the joint
null distribution to be invariant under these sign flips, as would hold for
independent seed contrasts symmetric about zero. The test concerns training
variation conditional on the observed cohort; it does not resample participants.
We applied Holm's step-down correction to the three candidate p-values
\citep{holm1979simple}, sorting by p-value and using the declared candidate
order only to break ties.

The prespecified rule required all five conditions: at least 50 external
participants; adjusted $p<0.05$; a bootstrap lower bound above zero; positive
differences for at least eight seeds; and a worst-seed difference no lower than
$-0.01$. The 50-person threshold was a minimum sample-size requirement, not a
power calculation. We retained this rule after seeing the results.

\subsection{Exploratory analyses}
\paragraph{Training-set size.}
We trained the baseline, rich-statistics head, and a residual MLP on nested HBN
subsets of 200, 400, and 800 participants, using the same ten seeds and fixed
validation set. The MLP has four hidden units and 2,569 parameters, close to
the rich head's 2,561. All 90 runs were evaluated on the same 75 MIPDB
participants. We estimated both the candidate advantage at each size and the
change in that advantage between sizes. The six cell comparisons and six
size contrasts formed separate exploratory sign-flip/Holm families.

\paragraph{Representation depth.}
The layer-wise analysis compared linear probes at layers $-4$, $-3$, and $-2$
with the final-layer baseline, using the same ten seeds and MIPDB participants.
It used the joint bootstrap and a separate three-comparison sign-flip/Holm
family. Both follow-ups were planned after the primary protocol and leave its
decision unchanged.

\paragraph{Transfer to ds006780.}
We evaluated the HBN-trained baseline and rich-statistics heads from the
200- and 800-subject training conditions on ds006780. Temporal filtering,
standardization, window length, aggregation, and checkpoint selection followed
the same protocol. The montage adapter instead selected 64 BioSemi EEG channels,
excluded auxiliary and trigger channels, used canonical electrode coordinates,
and resampled from 512 Hz. The source reference was undocumented and was
preserved rather than replaced.

The main comparison in this secondary analysis was the rich head versus the
baseline at 800 training subjects. Before adding external age labels to the
final cohort manifest, we specified a maximum 95\% interval width of
\DsPrecisionMaxWidth{}, corresponding to an intended precision of approximately
$\pm0.03$ Pearson units. We used the same 10,000-iteration joint bootstrap.
This precision requirement is distinct from evidence that an effect is positive.
